% GENERATED FILE. Edit canonical lessons, then run npm run generate:books.
% canonical-source-hash: fnv1a64:93cc62636cec184c
% canonical-lessons: SA-C539-apramattah, SA-C539-anavadhanam, SA-C539-jagarukah, SA-C539-asankate, SA-C539-samraksati

\chapter{Being Careful}
\label{ch:sa-being-careful}
\hlchaptermodality{539}

\begin{chapteropening}
\textbf{By the end of this chapter:} \emph{I can warn someone and say what happens if they are careless.}

Complete the last lesson of chapter 539: I can warn someone and say what happens if they are careless.
\end{chapteropening}

\section[\texorpdfstring{\sk{अप्रमत्तः} (apramattaḥ)}{अप्रमत्तः (apramattaḥ)}]{\sk{अप्रमत्तः} (apramattaḥ) — heedful}

\label{lesson:SA-C539-apramattah}



\begin{quote}
Before the new one: say the Sanskrit for an astrologer, then the Sanskrit for draws near.
\end{quote}

\subsection*{You'll want to know: \sk{अप्रमत्तः}}
\textbf{\sk{अप्रमत्तः}} — \emph{apramattaḥ} — \textquotedblleft{}heedful\textquotedblright{}.

\begin{quote}
\sk{यदि} \sk{त्वम्} \sk{अप्रमत्तः} \sk{असि}, \sk{तर्हि} \sk{त्वं} \sk{न} \sk{पतसि।}
\end{quote}

\begin{quote}
\emph{yadi tvam apramattaḥ asi, tarhi tvaṃ na patasi.}
\end{quote}

\begin{quote}
If you are heedful, you will not fall.
\end{quote}

\begin{cousinweb}
\textbf{\sk{अप्रमत्तः}} (\emph{apramattaḥ}) is built on \emph{a}, \textquotedblleft{}not\textquotedblright{}, and \emph{pramatta}, \textquotedblleft{}careless\textquotedblright{}, from \emph{pra} and the root √mad, \textquotedblleft{}to be drunk\textquotedblright{}; MW gives \emph{apramatta} as not careless, careful, attentive, vigilant (MW).
\end{cousinweb}

\begin{grammarlens}[title={What you've built}]
One more word for saying what happens if.
\end{grammarlens}

\subsection*{Your turn}
\begin{itemize}
\raggedright
  \item \emph{Say it:} \emph{apramattaḥ}
  \item \emph{Say it:} \emph{apramattaḥ}, once more
  \item \emph{Say it:} \emph{yadi tvam apramattaḥ asi, tarhi tvaṃ na patasi.} — If you are heedful, you will not fall.
  \item \emph{Recall:} say the Sanskrit for an astrologer, then the Sanskrit for draws near, then say \emph{apramattaḥ} again
\end{itemize}

\begin{tcolorbox}[breakable,colback=teal!4,colframe=teal!35!black,title={Before you move on}]
What does \sk{अप्रमत्तः} mean? (\textquotedblleft{}Heedful\textquotedblright{}.) Say it once more.
\end{tcolorbox}
\section[\texorpdfstring{\sk{अनवधानम्} (anavadhānam)}{अनवधानम् (anavadhānam)}]{\sk{अनवधानम्} (anavadhānam) — inattention}

\label{lesson:SA-C539-anavadhanam}



\begin{quote}
Before the new one: say the Sanskrit for draws near, then the Sanskrit for heedful.
\end{quote}

\subsection*{You'll want to know: \sk{अनवधानम्}}
\textbf{\sk{अनवधानम्}} — \emph{anavadhānam} — \textquotedblleft{}inattention\textquotedblright{}.

\begin{quote}
\sk{यदि} \sk{अनवधानं} \sk{भवति}, \sk{तर्हि} \sk{हानिः} \sk{अपि} \sk{भवति।}
\end{quote}

\begin{quote}
\emph{yadi anavadhānaṃ bhavati, tarhi hāniḥ api bhavati.}
\end{quote}

\begin{quote}
Where there is inattention, there is loss too.
\end{quote}

A neuter noun.

\begin{cousinweb}
\textbf{\sk{अनवधानम्}} (\emph{anavadhānam}) is built on \emph{an}, \textquotedblleft{}not\textquotedblright{}, and \emph{avadhāna}, \textquotedblleft{}attention\textquotedblright{}, from \emph{ava} and the root √dhā, \textquotedblleft{}to put\textquotedblright{}; MW gives \emph{anavadhāna} as inattention, inadvertence (MW).
\end{cousinweb}

\begin{grammarlens}[title={What you've built}]
One more word for saying what happens if.
\end{grammarlens}

\subsection*{Your turn}
\begin{itemize}
\raggedright
  \item \emph{Say it:} \emph{anavadhānam}
  \item \emph{Say it:} \emph{anavadhānam}, once more
  \item \emph{Say it:} \emph{yadi anavadhānaṃ bhavati, tarhi hāniḥ api bhavati.} — Where there is inattention, there is loss too.
  \item \emph{Recall:} say the Sanskrit for draws near, then the Sanskrit for heedful, then say \emph{anavadhānam} again
\end{itemize}

\begin{tcolorbox}[breakable,colback=teal!4,colframe=teal!35!black,title={Before you move on}]
What does \sk{अनवधानम्} mean? (\textquotedblleft{}Inattention\textquotedblright{}.) Say it once more.
\end{tcolorbox}
\section[\texorpdfstring{\sk{जागरूकः} (jāgarūkaḥ)}{जागरूकः (jāgarūkaḥ)}]{\sk{जागरूकः} (jāgarūkaḥ) — watchful}

\label{lesson:SA-C539-jagarukah}



\begin{quote}
Before the new one: say the Sanskrit for heedful, then the Sanskrit for inattention.
\end{quote}

\subsection*{You'll want to know: \sk{जागरूकः}}
\textbf{\sk{जागरूकः}} — \emph{jāgarūkaḥ} — \textquotedblleft{}watchful\textquotedblright{}.

\begin{quote}
\sk{यदि} \sk{रक्षकः} \sk{जागरूकः} \sk{अस्ति}, \sk{तर्हि} \sk{चौरः} \sk{न} \sk{आगच्छति।}
\end{quote}

\begin{quote}
\emph{yadi rakṣakaḥ jāgarūkaḥ asti, tarhi cauraḥ na āgacchati.}
\end{quote}

\begin{quote}
If the guard is watchful, the thief does not come.
\end{quote}

\begin{cousinweb}
\textbf{\sk{जागरूकः}} (\emph{jāgarūkaḥ}) comes from the root √jāgṛ, \textquotedblleft{}to be awake\textquotedblright{}, with the ending \emph{-ūka}; MW gives \emph{jāgarūka} as wakeful, watchful (MW).
\end{cousinweb}

\begin{grammarlens}[title={What you've built}]
One more word for saying what happens if.
\end{grammarlens}

\subsection*{Your turn}
\begin{itemize}
\raggedright
  \item \emph{Say it:} \emph{jāgarūkaḥ}
  \item \emph{Say it:} \emph{jāgarūkaḥ}, once more
  \item \emph{Say it:} \emph{yadi rakṣakaḥ jāgarūkaḥ asti, tarhi cauraḥ na āgacchati.} — If the guard is watchful, the thief does not come.
  \item \emph{Recall:} say the Sanskrit for heedful, then the Sanskrit for inattention, then say \emph{jāgarūkaḥ} again
\end{itemize}

\begin{tcolorbox}[breakable,colback=teal!4,colframe=teal!35!black,title={Before you move on}]
What does \sk{जागरूकः} mean? (\textquotedblleft{}Watchful\textquotedblright{}.) Say it once more.
\end{tcolorbox}
\section[\texorpdfstring{\sk{आशंकते} (āśaṅkate)}{आशंकते (āśaṅkate)}]{\sk{आशंकते} (āśaṅkate) — suspects; fears}

\label{lesson:SA-C539-asankate}



\begin{quote}
Before the new one: say the Sanskrit for inattention, then the Sanskrit for watchful.
\end{quote}

\subsection*{You'll want to know: \sk{आशंकते}}
\textbf{\sk{आशंकते}} — \emph{āśaṅkate} — \textquotedblleft{}suspects; fears\textquotedblright{}.

\begin{quote}
\sk{यदि} \sk{माता} \sk{चौरम्} \sk{आशंकते}, \sk{तर्हि} \sk{सा} \sk{द्वारं} \sk{पिदधाति।}
\end{quote}

\begin{quote}
\emph{yadi mātā cauram āśaṅkate, tarhi sā dvāraṃ pidadhāti.}
\end{quote}

\begin{quote}
If my mother suspects a thief, she shuts the door.
\end{quote}

Past (she …): \sk{सा} \sk{चौरम्} \sk{आशंकितवती}, \emph{sā cauram āśaṅkitavatī}.

\begin{cousinweb}
\textbf{\sk{आशंकते}} (\emph{āśaṅkate}) is built on \emph{ā} and the root √śaṅk, \textquotedblleft{}to doubt, to fear\textquotedblright{}; MW gives \emph{āśaṅk} as to suspect, to fear, to doubt (MW).
\end{cousinweb}

\begin{grammarlens}[title={What you've built}]
One more word for saying what happens if.
\end{grammarlens}

\subsection*{Your turn}
\begin{itemize}
\raggedright
  \item \emph{Say it:} \emph{āśaṅkate}
  \item \emph{Say it:} \emph{āśaṅkate}, once more
  \item \emph{Say it:} \emph{yadi mātā cauram āśaṅkate, tarhi sā dvāraṃ pidadhāti.} — If my mother suspects a thief, she shuts the door.
  \item \emph{Recall:} say the Sanskrit for inattention, then the Sanskrit for watchful, then say \emph{āśaṅkate} again
\end{itemize}

\begin{tcolorbox}[breakable,colback=teal!4,colframe=teal!35!black,title={Before you move on}]
What does \sk{आशंकते} mean? (\textquotedblleft{}Suspects; fears\textquotedblright{}.) Say it once more.
\end{tcolorbox}
\section[\texorpdfstring{\sk{संरक्षति} (saṃrakṣati)}{संरक्षति (saṃrakṣati)}]{\sk{संरक्षति} (saṃrakṣati) — keeps safe}

\label{lesson:SA-C539-samraksati}



\begin{quote}
Before the new one: say the Sanskrit for watchful, then the Sanskrit for suspects.
\end{quote}

\subsection*{You'll want to know: \sk{संरक्षति}}
\textbf{\sk{संरक्षति}} — \emph{saṃrakṣati} — \textquotedblleft{}keeps safe\textquotedblright{}.

\begin{quote}
\sk{यदि} \sk{त्वं} \sk{धनं} \sk{संरक्षसि}, \sk{तर्हि} \sk{तत्} \sk{न} \sk{नश्यति।}
\end{quote}

\begin{quote}
\emph{yadi tvaṃ dhanaṃ saṃrakṣasi, tarhi tat na naśyati.}
\end{quote}

\begin{quote}
If you keep the money safe, it will not be lost.
\end{quote}

Past (he …): \sk{सः} \sk{धनं} \sk{संरक्षितवान्}, \emph{saḥ dhanaṃ saṃrakṣitavān}.

\begin{cousinweb}
\textbf{\sk{संरक्षति}} (\emph{saṃrakṣati}) is built on \emph{sam}, \textquotedblleft{}wholly\textquotedblright{}, and the root √rakṣ, \textquotedblleft{}to protect\textquotedblright{}, which also gives \emph{rakṣati}; MW gives \emph{saṃrakṣ} as to protect, to guard, to keep, to secure (MW).
\end{cousinweb}

\begin{grammarlens}[title={What you've built}]
One more word for saying what happens if.
\end{grammarlens}

\subsection*{Your turn}
\begin{itemize}
\raggedright
  \item \emph{Say it:} \emph{saṃrakṣati}
  \item \emph{Say it:} \emph{saṃrakṣati}, once more
  \item \emph{Say it:} \emph{yadi tvaṃ dhanaṃ saṃrakṣasi, tarhi tat na naśyati.} — If you keep the money safe, it will not be lost.
  \item \emph{Recall:} say the Sanskrit for watchful, then the Sanskrit for suspects, then say \emph{saṃrakṣati} again
\end{itemize}

\begin{tcolorbox}[breakable,colback=teal!4,colframe=teal!35!black,title={Before you move on}]
What does \sk{संरक्षति} mean? (\textquotedblleft{}Keeps safe\textquotedblright{}.) Say it once more.
\end{tcolorbox}
